% ============================================================
%  SECTION 6: RESULTS
% ============================================================

\section{Results}
\label{sec:results}

\subsection{Main Effects}
\label{sec:results_main}

Figures~\ref{fig:synth_trends_par_q1}
and~\ref{fig:synth_trends_par_rank} provide visual evidence of the
compositional shift. The synthetic control, constructed from a weighted
combination of non-merit states
\parencite{abadie2010synthetic}, matches Tennessee almost exactly
through the pre-treatment period (1980--1984) for both outcomes. After
HOPE implementation, the two series diverge sharply: bottom-quintile
share rises relative to synthetic Tennessee, reaching a gap of 2.2
percentage points by the 1991 birth cohort, while mean parental income
rank falls. The pre-treatment fit is excellent for both outcomes, with
near-zero root mean squared prediction error.

\begin{figure}[t]
  \centering
  \includegraphics[width=0.85\textwidth]{fig_synth_trends_par_q1.pdf}
  \caption{Actual vs.\ Synthetic Tennessee: Bottom-Quintile Share at
    Public Colleges}
  \label{fig:synth_trends_par_q1}
  \vspace{1mm}
  \begin{minipage}{0.85\textwidth}
    \footnotesize
    \textit{Notes:} Synthetic Tennessee is constructed from a weighted
    combination of 22 non-merit states using the method of
    \textcite{abadie2010synthetic}, matching on pre-treatment
    bottom-quintile share (1980--1984). The vertical dashed line
    indicates the first HOPE-eligible birth cohort. The two series track
    closely in the pre-treatment period, then diverge after HOPE
    implementation. Data: \textcite{chetty2020mobility}.
  \end{minipage}
\end{figure}

\begin{figure}[t]
  \centering
  \includegraphics[width=0.85\textwidth]{fig_synth_trends_par_rank.pdf}
  \caption{Actual vs.\ Synthetic Tennessee: Mean Parental Income Rank at
    Public Colleges}
  \label{fig:synth_trends_par_rank}
  \vspace{1mm}
  \begin{minipage}{0.85\textwidth}
    \footnotesize
    \textit{Notes:} Synthetic Tennessee constructed as in
    Figure~\ref{fig:synth_trends_par_q1}, matching on pre-treatment
    mean parental income rank (1980--1984). Tennessee's actual parental
    income rank diverges downward after HOPE implementation, consistent
    with the progressive compositional shift documented in the
    bottom-quintile share results. Data: \textcite{chetty2020mobility}.
  \end{minipage}
\end{figure}

Table~\ref{tab:cs_main} reports the Callaway--Sant'Anna aggregate ATT
estimates. Following HOPE implementation, the share of bottom-quintile
students at Tennessee public colleges increased by 1.27 percentage points (SE =
0.0026, $p < 0.001$) and decreased mean parental income rank by 1.33
percentage points (SE = 0.0020, $p < 0.001$). Both effects are precisely estimated
and economically meaningful: the bottom-quintile share increase
represents a 9.2\% gain relative to the pre-treatment mean of 13.8\%.
In levels, the effect corresponds to approximately 11.2 additional
bottom-quintile students per cohort at each Tennessee public college and
34.4 fewer top-quintile students per cohort
(Section~\ref{sec:results_mechanisms}). Across 26 treated public
colleges, this implies roughly 290 additional low-income students and
890 fewer high-income students per birth cohort gaining access to or
exiting the public system. The effect grows with exposure horizon:
by $k = +5$ (the 1991 birth cohort), the treatment effect reaches
$+3.1$ percentage points, a 22\% gain over baseline. This growing path
reflects intensifying program impact on successive cohorts rather than
mechanical accumulation, since each cohort's composition is measured
independently (Section~\ref{sec:data}).

For comparison, the static TWFE estimator yields slightly larger point
estimates---1.54 pp for bottom-quintile share and $-1.85$ pp for
mean parental income rank (Table~\ref{tab:twfe_did})---consistent with the
known upward bias of TWFE in the presence of dynamic treatment effects.
The qualitative conclusions are identical across estimators.

These share changes reflect real compositional sorting rather than an
enrollment artifact. We confirm in Section~\ref{sec:results_mechanisms}
that the number of bottom-quintile students per cohort increases by 11.2
($p = 0.014$), the number of top-quintile students per cohort decreases
by 34.4 ($p < 0.001$), and total enrollment is unchanged ($-2.1$,
$p = 0.93$). The public sector absorbs more lower-income students while
losing upper-income students, with the total headcount held constant.

\subsection{Event Study Dynamics}
\label{sec:results_dynamics}

Figure~\ref{fig:four_estimators} presents the event-study estimates
across four heterogeneity-robust estimators
(Section~\ref{sec:robustness_estimators}). Focusing on the
Callaway--Sant'Anna path, for bottom-quintile share the
treatment effect is near zero at $k = 0$ and grows monotonically to
approximately 2.5 percentage points by $k = 5$. For parental income rank,
the trajectory is symmetric: a gradual decline reaching $-2.6$ points
by $k = 5$.

The growing post-treatment path is consistent with an intensifying
compositional shift across successive cohorts. Because each birth
cohort's composition is measured independently in the
\textcite{chetty2020mobility} data, the increasing treatment effect
reflects deeper program impact on later cohorts---through rising
take-up, greater awareness, and institutional adaptation---rather than
mechanical accumulation of treated students in a fixed enrollment stock.

Pre-treatment coefficients under the bordering-state control group are
centered around zero with no systematic trend
(Figure~\ref{fig:four_estimators}). Under the national control group,
pre-treatment coefficients at $k = -5$ and $k = -4$ show some
divergence, as discussed in Section~\ref{sec:empirical_strategy}.

\subsection{Quintile Decomposition}
\label{sec:results_quintile}

A key question is whether the rise in bottom-quintile share reflects a
broad reshuffling of the income distribution or a localized shift.
Table~\ref{tab:quintile_decomp} decomposes the composition effect
across all five income quintiles using the TWFE static specification.

The results reveal a strikingly monotonic gradient:

\begin{center}
\begin{tabular}{lccl}
\toprule
Quintile & Coefficient & SE & \\
\midrule
Q1 (bottom) & $+0.0154$ & $(0.0030)$ & $^{***}$ \\
Q2 & $+0.0114$ & $(0.0025)$ & $^{***}$ \\
Q3 & $+0.0047$ & $(0.0018)$ & $^{*}$ \\
Q4 & $-0.0136$ & $(0.0021)$ & $^{***}$ \\
Q5 (top) & $-0.0179$ & $(0.0025)$ & $^{***}$ \\
\bottomrule
\end{tabular}
\end{center}

Every quintile's coefficient has the predicted sign, and the gradient
is monotonically ordered: $|\hat\beta_1| > |\hat\beta_2| >
|\hat\beta_3|$ and $|\hat\beta_5| > |\hat\beta_4|$. The bottom three
quintiles gain share at public colleges; the top two lose share.
Statistical significance is strongest at the tails (Q1 and Q5) and
weakest in the middle (Q2 and Q3). This pattern is expected for two
reasons. First, the theoretical mechanism generates the largest effects
at the extremes: the proportional price reduction $S/y_i$ is largest
for the lowest-income families and the substitution incentive is
strongest for the highest-income families, while middle-quintile
families experience smaller behavioral responses in both directions.
Second, the tails of the income distribution exhibit greater
cross-college variance---bottom- and top-quintile shares vary more
across institutions than middle-quintile shares, which cluster near
0.20---giving the estimator more identifying variation. The combination
of larger true effects and greater statistical power at the extremes
means that Q1 and Q5 are precisely estimated, Q4 is significant, and
Q2--Q3 are directionally consistent but noisier.\footnote{Under the
Callaway--Sant'Anna estimator, Q2 is significant ($p < 0.001$) but Q3
is not ($p = 0.38$). Under TWFE, all five are significant. Under the
bordering-state specification with three estimators
(Figure~\ref{fig:quintile_event_studies}, Appendix), Q1 and Q5 are
significant across all estimators, Q4 is significant under TWFE and
Sun--Abraham, and Q2--Q3 are directionally consistent. The monotonic
ordering of point estimates is preserved across all specifications.}

\begin{figure}[t]
  \centering
  \includegraphics[width=0.75\textwidth]{fig_tn_quintile_gradient.pdf}
  \caption{Quintile Gradient of Compositional Change Following
    Tennessee HOPE}
  \label{fig:quintile_gradient}
  \vspace{1mm}
  \begin{minipage}{0.75\textwidth}
    \footnotesize
    \textit{Notes:} \textcite{callaway2021difference} ATT estimates for
    each parental income quintile share. The effect is monotonically
    ordered: bottom-quintile students gain the most representation and
    top-quintile students lose the most. Significance is strongest at the
    tails (Q1, Q5) where both the theoretical effect and cross-college
    variance are largest. Bars denote 95\% confidence intervals.
    Data: \textcite{chetty2020mobility}.
  \end{minipage}
\end{figure}

Figure~\ref{fig:quintile_gradient} displays the quintile gradient
graphically. The coefficients descend smoothly from Q1 ($+1.3$ pp)
through Q3 ($+0.2$ pp) to Q5 ($-1.1$ pp). The quintile-specific
event studies (Figure~\ref{fig:quintile_event_studies}, Appendix)
confirm that all five quintiles exhibit gradual divergence after
treatment onset, with three estimators agreeing on the direction and
dynamics at each quintile.

The quintile shares sum to one by construction. The finding that Q1
gains the most share and Q5 loses the most share implies that the
composition shift is driven by reallocation between the extremes of
the income distribution rather than by changes near the median. The
monotonic gradient also serves as a discriminating test of the HOPE
channel: a generic Tennessee-specific shock---whether driven by state
fiscal conditions, demographic trends, or concurrent reforms---would be
unlikely to produce five coefficients that are perfectly ordered in
magnitude. This specific pattern is a unique prediction of the
price-reduction mechanism in the three-destination model, where the
proportional cost reduction $S/y_i$ is monotonically decreasing in
income.

\subsection{Mechanism Evidence}
\label{sec:results_mechanisms}

\subsubsection{Level Decomposition: Shares vs.\ Counts}

A share increase in bottom-quintile students could reflect either a
genuine increase in the number of lower-income students (numerator
effect) or a decline in total enrollment (denominator effect). We
address this ecological inference concern by estimating effects on
count outcomes: Q1 count (number of bottom-quintile students per
cohort), Q5 count, and total enrollment.\footnote{Count outcomes are
constructed as quintile share $\times$ cohort enrollment. Because
enrollment varies substantially across colleges (coefficient of
variation $\approx 1.0$), count-based estimates are considerably
noisier than share-based estimates. The critical test is whether
\emph{total} enrollment changes, which is well-powered because it does
not inherit share-level noise. The directional signs of Q1 and Q5 counts
provide supporting evidence but should be interpreted with appropriate
caution given the wider confidence intervals.}

The critical finding is that total enrollment at Tennessee public
colleges is essentially unchanged after HOPE ($-2.1$ students per
cohort, $p = 0.93$). This rules out the denominator-artifact
explanation: the share changes cannot be driven by enrollment
contraction. The directional evidence is consistent with real
compositional sorting: Q1 count increases by 11.2 per cohort ($p =
0.014$) and Q5 count decreases by 34.4 per cohort ($p < 0.001$),
indicating that more bottom-quintile students enter the public sector
while more top-quintile students exit, with the total headcount held
constant (Table~\ref{tab:level_decomp}).

\subsubsection{Tier Heterogeneity}

Table~\ref{tab:tier_het} reports separate estimates for two-year and
four-year institutions.

At two-year colleges, the bottom-quintile share increases by 1.81
percentage points (SE = 0.0034, $p < 0.001$) and mean parental income
rank falls by 1.73 percentage points. These institutions show a loss of 22.9 Q5
students per cohort ($p = 0.001$) but no significant change in total
enrollment. The evidence is consistent with \emph{displacement}: lower-income
students replace higher-income students at community colleges, with the
overall seat count unchanged.

At four-year colleges, the bottom-quintile share increases by a smaller
0.70 percentage points (SE = 0.0024, $p = 0.008$), but total enrollment
increases by 138.5 students per cohort ($p < 0.001$). Four-year
institutions gain 19.6 Q1 students and lose 39.2 Q5 students. The evidence is consistent with \emph{expansion}: these institutions
grow their enrollment while simultaneously experiencing downward
compositional sorting.

The two tiers thus exhibit different adjustment channels, consistent
with Prediction P5. Community colleges adjust on the intensive margin
(who fills existing seats), while four-year institutions adjust on
the extensive margin (adding seats while changing composition).

\subsubsection{Tennessee Privates: No Spillover}

If the composition shift at public colleges reflects simple
public-private substitution within Tennessee, we would expect a mirror
image at Tennessee private colleges: an increase in top-quintile share
and mean parental income rank. We test this directly.

Table~\ref{tab:spillover} reports Callaway--Sant'Anna ATTs comparing
Tennessee private colleges to private colleges in control states. The
effect on bottom-quintile share is 0.11 percentage points (SE = 0.0011,
$p = 0.30$)---small, positive, and statistically insignificant. The
effect on mean parental income rank is $-0.46$ percentage points (SE = 0.0018, $p =
0.016$), marginally significant but in the same direction as the
public-sector effect rather than the opposite direction.

Tennessee private colleges do not absorb the top-quintile students who
leave the public sector. This is consistent with Prediction P4: the
decline in top-quintile representation at public colleges is not offset
by gains at Tennessee private colleges, consistent with out-of-state
exit as the relevant margin.
Figure~\ref{fig:pub_vs_priv} displays the contrast visually:
the public-sector event study shows a clear progressive shift while
the private-sector event study is flat, confirming that the
compositional change is specific to the public sector.
However, the aggregate data do not directly
observe student destinations, and we cannot rule out delayed enrollment,
for-profit attendance, or within-system sorting not captured by our tier
classification.

\begin{figure}[H]
  \centering
  \includegraphics[width=\textwidth]{fig_public_vs_private_2x2_combined.pdf}
  \caption{Public vs.\ Private Colleges: Event Studies for Both Outcomes}
  \label{fig:pub_vs_priv}
  \vspace{1mm}
  \begin{minipage}{\textwidth}
    \footnotesize
    \textit{Notes:} Left column: four heterogeneity-robust estimators
    for Tennessee public colleges vs.\ public colleges in 26 non-merit
    states. Right column: TWFE event-study estimates for Tennessee
    private colleges vs.\ private colleges in control states. Top row:
    bottom-quintile share. Bottom row: mean parental income rank.
    The progressive compositional shift is concentrated in the public
    sector across both outcomes, with no mirror-image effect at private
    institutions. Vertical bars denote 95\% confidence intervals.
    Data: \textcite{chetty2020mobility}.
  \end{minipage}
\end{figure}

\subsubsection{Where Did Top-Quintile Students Go?}

The tier heterogeneity results (Table~\ref{tab:tier_het}) shed light on
the Q5 exit pattern. Two-year colleges lost 22.9 Q5 students per cohort
($p = 0.001$) with no significant enrollment change, while four-year
colleges lost 39.2 Q5 students ($p = 0.034$) alongside enrollment
\emph{growth} of 138.5 students ($p < 0.001$). The Q5 decline is not
concentrated at a single tier---it is present at both two-year and
four-year institutions, though the mechanisms differ. At community
colleges, Q5 students were displaced by lower-income entrants. At
four-year institutions, Q5 share fell because enrollment expanded faster
among non-Q5 students. In neither case do the displaced Q5 students appear at Tennessee private
colleges, reinforcing the interpretation that the relevant exit margin
is outside Tennessee's private sector.


\subsubsection{Within-Tennessee Public vs.\ Private Comparison}

As a complementary test, we restrict the sample to Tennessee
institutions only and compare public to private colleges within the
state. This eliminates any concerns about cross-state confounds.

The within-Tennessee DiD estimate for bottom-quintile share is 1.42
percentage points (SE = 0.0037, $p < 0.001$) and for mean parental income rank
is $-2.14$ percentage points (SE = 0.0032, $p < 0.001$). These estimates are
slightly larger than the cross-state estimates, consistent with
Tennessee private colleges serving as a clean within-state control.
The sample is small ($N = 571$), but the effects are precisely estimated
and confirm that the composition shift is specific to the public sector.

\subsection{Spatial Heterogeneity}
\label{sec:spatial}

The three-destination model predicts that the composition shift varies
with institutional geography (P6). Colleges closer to state borders
face stronger outside-option competition---the kind of cross-market
pressure emphasized by \textcite{hoxby2000peer}---which should amplify
the sorting mechanism. We test this prediction by estimating heterogeneous
effects along two spatial dimensions: border proximity and region.

\textbf{Border vs.\ interior.} Institutions in Tennessee counties that
border another state show a 38\% larger progressive shift than interior
institutions: the bottom-quintile share ATT is $+1.89$ pp (SE $=
0.003$, $p < 0.001$) at border colleges versus $+1.37$ pp (SE $=
0.003$, $p < 0.001$) at interior colleges. The monotonic quintile
gradient holds in both subsamples, but the gradient is steeper at
border colleges. This is consistent with P6: where students have
accessible out-of-state alternatives, HOPE's price reduction has a
larger behavioral effect on the margin between public enrollment and
exit.

\textbf{Regional variation.} The effect varies strikingly across
Tennessee's three Grand Divisions (Table~\ref{tab:spatial}). West
Tennessee shows the largest shift ($+3.13$ pp, $p < 0.001$), roughly
triple the magnitude in East Tennessee ($+1.13$ pp, $p < 0.001$) and
Middle Tennessee ($+1.03$ pp, $p = 0.002$). The monotonic quintile
gradient holds in all three regions. The concentration in West
Tennessee---where the Memphis metropolitan area's proximity to
Mississippi and Arkansas creates strong cross-border competition, and a
high density of community colleges favors both access expansion and
competitive retention---is consistent with the model's spatial
predictions.\footnote{West Tennessee contains only 6 of the 26 public
colleges. The large point estimate should be interpreted with this
small-sample caveat in mind.}

Figure~\ref{fig:spatial_map} displays the geographic distribution of
the effect. Each of Tennessee's 26 public colleges is marked at its
county location, with the county fill reflecting the regional ATT
magnitude. The concentration of the effect in West Tennessee---and
particularly in border counties adjacent to Mississippi and
Arkansas---is visually striking.

\begin{figure}[H]
  \centering
  \includegraphics[width=\textwidth]{fig_spatial_heterogeneity.pdf}
  \caption{Geographic Distribution of HOPE's Compositional Effect}
  \label{fig:spatial_map}
  \vspace{1mm}
  \begin{minipage}{\textwidth}
    \footnotesize
    \textit{Notes:} County fill reflects the regional bottom-quintile
    share ATT (West: $+3.1$ pp, East: $+1.1$ pp, Middle: $+1.0$ pp),
    with border counties receiving a 38\% boost per the border vs.\
    interior heterogeneity estimate. Navy markers indicate public
    college locations: circles = four-year universities, triangles =
    two-year colleges. Border counties are outlined. Neighboring states
    labeled to illustrate out-of-state exit destinations.
    Data: \textcite{chetty2020mobility}.
  \end{minipage}
\end{figure}

\begin{table}[t]
\centering
\begin{threeparttable}
\caption{Spatial Heterogeneity in HOPE's Compositional Effect}
\label{tab:spatial}
\begin{tabular}{lccc}
\toprule
 & Q1 Share ATT & SE & Colleges \\
\midrule
\textit{By Border Status} & & & \\
\quad Border counties & $0.0189^{***}$ & $(0.003)$ & 8 \\
\quad Interior counties & $0.0137^{***}$ & $(0.003)$ & 18 \\
\midrule
\textit{By Region} & & & \\
\quad East Tennessee & $0.0113^{***}$ & $(0.003)$ & 9 \\
\quad Middle Tennessee & $0.0103^{**}$ & $(0.003)$ & 11 \\
\quad West Tennessee & $0.0313^{***}$ & $(0.003)$ & 6 \\
\bottomrule
\end{tabular}
\begin{tablenotes}[flushleft]
\footnotesize
\item \textit{Notes:} TWFE static difference-in-differences estimates.
  College and cohort fixed effects included. Standard errors clustered
  at the state level in parentheses. Border counties share a boundary
  with another state. Regional definitions follow Tennessee's Grand
  Divisions. $^{***}p<0.01$, $^{**}p<0.05$, $^{*}p<0.10$.
  Data: \textcite{chetty2020mobility}.
\end{tablenotes}
\end{threeparttable}
\end{table}

\subsection{Summary of Theoretical Predictions}

Table~\ref{tab:predictions} maps each prediction from the
three-destination model to its empirical test and result.

\begin{table}[H]
\centering
\begin{threeparttable}
\caption{Model Predictions and Empirical Tests}
\label{tab:predictions}
\begin{tabular}{llll}
\toprule
Prediction & Test & Section & Result \\
\midrule
P1: Q1 share rises & C\&S ATT & \ref{sec:results_main} & $+1.27$ pp ($p < 0.001$) \\
P2: Monotonic gradient & Quintile decomp. & \ref{sec:results_quintile} & 5/5 ordered \\
P3: Sorting, not expansion & Level decomp. & \ref{sec:results_mechanisms} & Total $\Delta = -2.1$ ($p = 0.93$) \\
P4: No mirror at privates & Spillover test & \ref{sec:results_mechanisms} & Null ($p = 0.30$) \\
P5: Tier heterogeneity & 2-yr vs.\ 4-yr & \ref{sec:results_mechanisms} & Different channels \\
P6: Spatial gradient & Border vs.\ interior & \ref{sec:spatial} & Border 38\% larger \\
P7: Cross-state variation & TN, GA, WV & \ref{sec:discussion} & Design predicts direction \\
\bottomrule
\end{tabular}
\begin{tablenotes}[flushleft]
\footnotesize
\item \textit{Notes:} Predictions are derived in Section~\ref{sec:theory}.
  All Tennessee results use the preferred treatment timing ($t^* = 1986$).
  P7 is evaluated in the Discussion using cross-state evidence.
\end{tablenotes}
\end{threeparttable}
\end{table}
